\chapter{Operating Systems Laboratory Projects and Practical Engineering}
\label{chap:labs}

A fundamental principle established by Andrew S. Tanenbaum in \textit{Operating Systems: Design and Implementation} is that an operating system cannot be truly mastered solely through theoretical lectures and mathematical queuing proofs. True mastery requires getting one's hands dirty: modifying the kernel, implementing new system calls, debugging hardware page faults, tuning scheduling algorithms, and writing real device drivers that communicate with physical or emulated hardware.

This chapter presents eight comprehensive, university-level laboratory projects designed for students, educators, and systems software engineers studying Hydra. Each project includes:
\begin{enumerate}
    \item \textbf{Problem Statement and Architectural Motivation}: Why this feature is needed in modern operating systems.
    \item \textbf{Design Specification}: Precise API signatures, data structures, and state transitions.
    \item \textbf{Implementation Guide}: Concrete, runnable Rust and C source code snippets illustrating the kernel modifications.
    \item \textbf{Verification Test Harness}: User-space test programs that validate correct behavior and edge-case handling.
    \item \textbf{Grading Rubric and Advanced Challenges}: Quantitative criteria for evaluation and exploratory extensions.
\end{enumerate}

% ==============================================================================
\section{Lab 1: Process CPU Time Accounting and \texttt{sys\_get\_cpu\_usage}}
\label{sec:lab1}
% ==============================================================================

\subsection{Architectural Motivation}

Operating system task managers (such as \texttt{top}, \texttt{htop}, and the KDE Plasma System Monitor) require accurate metrics detailing how much CPU time each active process has consumed. In modern multicore processors, measuring time solely by coarse timer ticks (e.g., $10\,\text{ms}$ PIT or HPET intervals) leads to severe quantization error: a process that executes for $1\,\text{ms}$ and yields before the timer interrupt fires might be billed for zero ticks, while a process that happens to be running during the interrupt is billed for the entire interval.

In this lab, students extend the Hydra microkernel to track CPU cycle consumption using the x86\_64 Time Stamp Counter (\texttt{RDTSC}) and expose this metric via a new microkernel system call: \texttt{sys\_get\_cpu\_usage}.

\subsection{Design Specification}

The system call is assigned number 105 in Hydra's native microkernel dispatch table:
\begin{lstlisting}[language=C]
struct CpuUsage {
    uint64_t user_cycles;    // CPU cycles spent in Ring 3 user mode
    uint64_t system_cycles;  // CPU cycles spent in Ring 0 kernel mode
    uint64_t voluntary_switches;   // Voluntary yields (e.g., IPC, sleep)
    uint64_t involuntary_switches; // Preemptions by timer interrupt
};

int sys_get_cpu_usage(uint64_t pid, struct CpuUsage *usage);
\end{lstlisting}

\subsection{Implementation Walkthrough}

First, we extend the Task Control Block (\texttt{src/task/tcb.rs}) to store cycle counters:

\begin{lstlisting}[language=Rust, caption={TCB Accounting Fields (\texttt{src/task/tcb.rs})}]
pub struct TaskControlBlock {
    // ... existing fields ...
    pub user_cycles: u64,
    pub system_cycles: u64,
    pub voluntary_switches: u64,
    pub involuntary_switches: u64,
    pub last_switch_tsc: u64,
}
\end{lstlisting}

Next, in the scheduler context-switch routine (\texttt{src/task/scheduler.rs}), we sample the hardware Time Stamp Counter using the \texttt{rdtsc} instruction and update the outgoing and incoming tasks:

\begin{lstlisting}[language=Rust, caption={Context Switch TSC Accounting (\texttt{src/task/scheduler.rs})}]
pub fn switch_to(&mut self, next_idx: usize, involuntary: bool) {
    let now_tsc = unsafe { core::arch::x86_64::_rdtsc() };
    
    // 1. Account elapsed cycles to the outgoing task
    if let Some(prev_idx) = self.current_task_idx {
        let prev_task = &mut self.tasks[prev_idx];
        let elapsed = now_tsc.saturating_sub(prev_task.last_switch_tsc);
        
        // For simplicity, attribute execution between switches
        prev_task.user_cycles += elapsed;
        if involuntary {
            prev_task.involuntary_switches += 1;
        } else {
            prev_task.voluntary_switches += 1;
        }
    }

    // 2. Prepare incoming task
    let next_task = &mut self.tasks[next_idx];
    next_task.last_switch_tsc = now_tsc;
    self.current_task_idx = Some(next_idx);

    // Perform architectural register switch...
}
\end{lstlisting}

Finally, we implement the system call handler in \texttt{src/syscall/user.rs}:

\begin{lstlisting}[language=Rust, caption={System Call Handler (\texttt{src/syscall/user.rs})}]
#[repr(C)]
pub struct CpuUsage {
    pub user_cycles: u64,
    pub system_cycles: u64,
    pub voluntary_switches: u64,
    pub involuntary_switches: u64,
}

pub fn sys_get_cpu_usage(pid: u64, out_ptr: *mut CpuUsage) -> i64 {
    if out_ptr.is_null() {
        return -14; // EFAULT
    }

    let sched = SCHEDULER.lock();
    let target = sched.find_task(TaskId::new(pid));
    
    match target {
        Some(task) => {
            let lock = task.lock();
            let usage = CpuUsage {
                user_cycles: lock.user_cycles,
                system_cycles: lock.system_cycles,
                voluntary_switches: lock.voluntary_switches,
                involuntary_switches: lock.involuntary_switches,
            };
            unsafe {
                core::ptr::write_volatile(out_ptr, usage);
            }
            0 // Success
        }
        None => -3, // ESRCH: No such process
    }
}
\end{lstlisting}

\subsection{Verification Test Harness}

\begin{lstlisting}[language=C, caption={Userland Test Program (\texttt{test\_cpu\_usage.c})}]
#include <stdio.h>
#include <stdint.h>
#include <unistd.h>
#include <assert.h>

struct CpuUsage {
    uint64_t user_cycles;
    uint64_t system_cycles;
    uint64_t voluntary_switches;
    uint64_t involuntary_switches;
};

int main() {
    struct CpuUsage u1, u2;
    uint64_t my_pid = getpid();

    int res = syscall(105, my_pid, &u1);
    assert(res == 0);

    // Burn some CPU cycles
    volatile uint64_t sum = 0;
    for (int i = 0; i < 10000000; i++) {
        sum += i;
    }

    res = syscall(105, my_pid, &u2);
    assert(res == 0);
    assert(u2.user_cycles > u1.user_cycles);

    printf("[PASS] Lab 1: CPU cycles measured: %lu -> %lu (delta: %lu cycles)\n",
           u1.user_cycles, u2.user_cycles, u2.user_cycles - u1.user_cycles);
    return 0;
}
\end{lstlisting}

% ==============================================================================
\section{Lab 2: Proportional-Share Stride Scheduling}
\label{sec:lab2}
% ==============================================================================

\subsection{Architectural Motivation}

Standard Round-Robin and FIFO schedulers treat all processes equally, giving every process an identical time slice. However, real-world systems run workloads with vastly different priorities: video rendering threads and audio playback daemons require guaranteed throughput and low latency, while background compilation tasks should only consume spare cycles.

In this lab, students replace Hydra's default Round-Robin scheduler with a deterministic \textbf{Stride Scheduler} (Waldspurger and Weihl, 1995). Unlike probabilistic lottery scheduling, stride scheduling guarantees exact proportional CPU allocation over finite intervals.

\subsection{Mathematical Derivation of Stride Scheduling}

Each task $i$ is assigned a number of tickets $T_i \ge 1$. The stride $S_i$ is inversely proportional to the ticket allocation:
\begin{equation}
S_i = \frac{L}{T_i}
\end{equation}
where $L$ is a large constant (e.g., $L = 10{,}000$). Each task maintains a \textit{pass} value $P_i$. At every scheduling decision, the scheduler selects the task with the minimum pass value:
\begin{equation}
\text{next} = \arg\min_{i} P_i
\end{equation}
Once task $i$ runs for one quantum, its pass value advances by its stride:
\begin{equation}
P_i \leftarrow P_i + S_i
\end{equation}

\subsection{Implementation Walkthrough}

\begin{lstlisting}[language=Rust, caption={Stride Scheduler Engine (\texttt{src/task/stride.rs})}]
pub struct StrideTask {
    pub id: TaskId,
    pub tickets: u32,
    pub stride: u64,
    pub pass: u64,
}

pub struct StrideScheduler {
    tasks: alloc::vec::Vec<StrideTask>,
    large_constant: u64,
}

impl StrideScheduler {
    pub fn new() -> Self {
        Self {
            tasks: alloc::vec::Vec::new(),
            large_constant: 10_000,
        }
    }

    pub fn add_task(&mut self, id: TaskId, tickets: u32) {
        let tickets = tickets.max(1);
        let stride = self.large_constant / tickets as u64;
        // New tasks initialize pass to the current minimum pass to avoid starvation
        let min_pass = self.tasks.iter().map(|t| t.pass).min().unwrap_or(0);
        self.tasks.push(StrideTask { id, tickets, stride, pass: min_pass });
    }

    pub fn pick_next(&mut self) -> Option<TaskId> {
        if self.tasks.is_empty() { return None; }

        // Find task with minimum pass value
        let mut min_idx = 0;
        let mut min_pass = self.tasks[0].pass;

        for (i, task) in self.tasks.iter().enumerate().skip(1) {
            if task.pass < min_pass {
                min_pass = task.pass;
                min_idx = i;
            }
        }

        // Advance pass value by stride
        self.tasks[min_idx].pass = self.tasks[min_idx].pass.wrapping_add(self.tasks[min_idx].stride);
        Some(self.tasks[min_idx].id)
    }
}
\end{lstlisting}

% ==============================================================================
\section{Lab 3: Kernel Object Cache (Slab Allocator)}
\label{sec:lab3}
% ==============================================================================

\subsection{Architectural Motivation}

Operating system kernels frequently allocate and free thousands of small, identical data structures per second (e.g., Task Control Blocks, File Descriptors, Socket Buffers, and VFS Inodes). Using a general-purpose heap allocator (such as a first-fit linked list or buddy allocator) causes severe external memory fragmentation and cache misses.

In this lab, students implement a \textbf{Slab Allocator} based on Jeff Bonwick's seminal SunOS paper (1994). The slab allocator groups identical objects into contiguous memory slabs carved from $4\,\text{KiB}$ pages, maintaining a singly-linked free list inside the free objects themselves.

\subsection{Implementation Walkthrough}

\begin{lstlisting}[language=Rust, caption={Slab Allocator Cache (\texttt{src/memory/slab.rs})}]
pub struct Cache {
    name: &'static str,
    object_size: usize,
    free_list: Option<*mut FreeNode>,
    pages_allocated: usize,
}

struct FreeNode {
    next: Option<*mut FreeNode>,
}

impl Cache {
    pub const fn new(name: &'static str, size: usize) -> Self {
        // Enforce alignment and minimum size for FreeNode pointer
        let object_size = if size < core::mem::size_of::<usize>() {
            core::mem::size_of::<usize>()
        } else {
            (size + 7) & !7 // 8-byte aligned
        };
        Self { name, object_size, free_list: None, pages_allocated: 0 }
    }

    /// Allocates a single pre-sized object from the cache.
    pub fn alloc(&mut self) -> *mut u8 {
        if let Some(node_ptr) = self.free_list {
            // Pop head of free list
            unsafe {
                self.free_list = (*node_ptr).next;
                node_ptr as *mut u8
            }
        } else {
            // Grow slab by allocating a new 4 KiB frame
            self.grow();
            self.alloc()
        }
    }

    /// Returns an object to the cache free list.
    pub fn free(&mut self, ptr: *mut u8) {
        if ptr.is_null() { return; }
        unsafe {
            let node = ptr as *mut FreeNode;
            (*node).next = self.free_list;
            self.free_list = Some(node);
        }
    }

    fn grow(&mut self) {
        let frame = crate::memory::frame_allocator::allocate_frame()
            .expect("Out of memory for slab growth");
        let page_virt = frame.start_address().as_u64() as *mut u8;
        let objects_per_page = 4096 / self.object_size;

        for i in 0..objects_per_page {
            let obj_ptr = unsafe { page_virt.add(i * self.object_size) };
            self.free(obj_ptr);
        }
        self.pages_allocated += 1;
    }
}
\end{lstlisting}

% ==============================================================================
\section{Lab 4: Zero-Copy Shared Memory Inter-Process Communication}
\label{sec:lab4}
% ==============================================================================

\subsection{Architectural Motivation}

Message-passing IPC on microkernels traditionally copies payload bytes from the sender's user space into kernel buffers, and then from kernel buffers into the receiver's address space. For high-bandwidth transfers (such as video framebuffers, network streaming, and audio mixing), this double-copy mechanism saturates system memory buses.

In this lab, students implement zero-copy shared memory windows (\texttt{sys\_shm\_create} and \texttt{sys\_shm\_map}). Two distinct processes map the exact same physical memory frames into their respective virtual address spaces, achieving multi-gigabyte per second transfer rates with zero CPU copy overhead.

\begin{lstlisting}[language=Rust, caption={Shared Memory Page Mapping (\texttt{src/ipc/shm.rs})}]
pub struct ShmRegion {
    pub id: u64,
    pub frames: alloc::vec::Vec<x86_64::structures::paging::PhysFrame>,
    pub size_pages: usize,
}

pub fn map_shm_into_task(
    shm: &ShmRegion,
    task_page_table: &mut crate::memory::paging::PageTable,
    target_virt: x86_64::VirtAddr,
) -> Result<(), ()> {
    use x86_64::structures::paging::PageTableFlags as Flags;
    let flags = Flags::PRESENT | Flags::WRITABLE | Flags::USER_ACCESSIBLE;

    for (i, frame) in shm.frames.iter().enumerate() {
        let page = x86_64::structures::paging::Page::from_start_address(
            target_virt + (i * 4096) as u64
        ).map_err(|_| ())?;

        task_page_table.map_page(page, *frame, flags)?;
    }
    Ok(())
}
\end{lstlisting}

% ==============================================================================
\section{Lab 5: Real-Time Block Compression (LZ4) in OpenZFS}
\label{sec:lab5}
% ==============================================================================

\subsection{Architectural Motivation}

Disk I/O is typically the slowest component in modern computing systems. By compressing storage blocks in RAM before writing them to physical disks, an operating system achieves two simultaneous benefits: it reduces physical storage consumption and effectively increases disk throughput (since fewer physical sectors must be transferred across SATA/NVMe buses).

In this lab, students integrate the high-throughput \textbf{LZ4 compression algorithm} into Hydra's OpenZFS Data Management Unit (DMU) pipeline (\texttt{src/fs/zfs/mod.rs}).

\subsection{Design Specification}

When writing a file block, the ZFS storage pipeline:
\begin{enumerate}
    \item Compresses the $128\,\text{KiB}$ plaintext block into a temporary buffer using byte-token LZ4.
    \item If the compressed size is at least 12.5\% smaller than the original block (saving at least one $16\,\text{KiB}$ allocation unit), the compressed data is written to disk and the Dnode's \texttt{dn\_flags} byte records \texttt{ZIO\_COMPRESS\_LZ4}.
    \item If the block is non-compressible (e.g., encrypted or already compressed media), the uncompressed block is stored verbatim, avoiding wasted decompression overhead.
\end{enumerate}

\begin{lstlisting}[language=Rust, caption={LZ4 Compression Hook in ZFS DMU Pipeline (\texttt{src/fs/zfs/compress.rs})}]
//! # OpenZFS Block Compression Pipeline Hook
//!
//! Provides real-time block compression for the SPA and DMU tiers.

pub const ZIO_COMPRESS_OFF: u8 = 0;
pub const ZIO_COMPRESS_LZ4: u8 = 15;

/// Compresses a 128 KiB ZFS block using LZ4.
///
/// # Arguments
/// * `src` - Plaintext data slice (typically 128 KiB).
/// * `dst` - Target buffer for compressed payload.
///
/// # Returns
/// Compressed byte length on success, or None if the block was incompressible.
pub fn zfs_compress_lz4(src: &[u8], dst: &mut [u8]) -> Option<usize> {
    if src.len() > 128 * 1024 || dst.len() < src.len() {
        return None;
    }

    let mut in_pos = 0;
    let mut out_pos = 0;
    let mut hash_table = [0u16; 4096];

    while in_pos + 4 < src.len() && out_pos + 8 < dst.len() {
        let sequence = u32::from_le_bytes([
            src[in_pos], src[in_pos + 1], src[in_pos + 2], src[in_pos + 3]
        ]);
        let hash = ((sequence.wrapping_mul(2654435761)) >> 20) as usize & 4095;
        let match_pos = hash_table[hash] as usize;
        hash_table[hash] = in_pos as u16;

        if match_pos < in_pos && in_pos - match_pos < 65535 
           && src[match_pos..match_pos + 4] == src[in_pos..in_pos + 4] {
            // Found duplicate sequence! Compute match length
            let mut match_len = 4;
            while in_pos + match_len < src.len() 
                  && src[match_pos + match_len] == src[in_pos + match_len] 
                  && match_len < 255 {
                match_len += 1;
            }

            // Emit LZ4 token: 4-bit literal length (0) | 4-bit match length
            let token = ((match_len - 4).min(15) as u8) & 0x0F;
            dst[out_pos] = token;
            out_pos += 1;

            // Emit 16-bit match offset (little-endian)
            let offset = (in_pos - match_pos) as u16;
            dst[out_pos] = (offset & 0xFF) as u8;
            dst[out_pos + 1] = ((offset >> 8) & 0xFF) as u8;
            out_pos += 2;

            in_pos += match_len;
        } else {
            // Emit literal byte
            dst[out_pos] = 0x10; // Token: 1 literal, 0 match
            dst[out_pos + 1] = src[in_pos];
            out_pos += 2;
            in_pos += 1;
        }
    }

    // Verify 12.5% compression threshold (at least 1/8 savings)
    if out_pos < src.len() - (src.len() / 8) {
        Some(out_pos)
    } else {
        None // Incompressible block
    }
}
\end{lstlisting}

% ==============================================================================
\section{Lab 6: Asynchronous VirtIO-Net Driver Implementation}
\label{sec:lab6}
% ==============================================================================

\subsection{Architectural Motivation}

VirtIO is the industry-standard paravirtualization framework used by QEMU, KVM, and cloud hypervisors. Unlike physical network adapters (which require complex register protocols and proprietary descriptor layouts), VirtIO devices use standardized \textbf{Virtqueues} consisting of three memory rings:
\begin{enumerate}
    \item \textbf{Descriptor Table}: Buffers containing packet payloads.
    \item \textbf{Available Ring}: Driver-to-device notifications listing buffers ready for consumption.
    \item \textbf{Used Ring}: Device-to-driver notifications listing buffers completed by the hypervisor.
\end{enumerate}

In this lab, students implement a complete Ring~3 VirtIO-Net driver and benchmark its packet reception latency against the legacy e1000 driver.

\begin{lstlisting}[language=Rust, caption={Virtqueue Memory Structures (\texttt{src/net/virtio.rs})}]
/// VirtIO 1.0 Standard Descriptor (16 bytes)
#[repr(C)]
#[derive(Clone, Copy, Default)]
pub struct VirtqDesc {
    pub addr: u64,
    pub len: u32,
    pub flags: u16,
    pub next: u16,
}

pub const VIRTQ_DESC_F_NEXT: u16     = 1;
pub const VIRTQ_DESC_F_WRITE: u16    = 2;
pub const VIRTQ_DESC_F_INDIRECT: u16 = 4;

/// Virtqueue Available Ring (Driver -> Device)
#[repr(C)]
pub struct VirtqAvail {
    pub flags: u16,
    pub idx: u16,
    pub ring: [u16; 256],
    pub used_event: u16,
}

/// Virtqueue Used Ring (Device -> Driver)
#[repr(C)]
pub struct VirtqUsed {
    pub flags: u16,
    pub idx: u16,
    pub ring: [VirtqUsedElem; 256],
    pub avail_event: u16,
}

#[repr(C)]
#[derive(Clone, Copy, Default)]
pub struct VirtqUsedElem {
    pub id: u32,
    pub len: u32,
}
\end{lstlisting}

% ==============================================================================
\section{Lab 7: Linux \texttt{io\_uring} Asynchronous Ring Buffer System Calls}
\label{sec:lab7}
% ==============================================================================

\subsection{Architectural Motivation}

Modern web servers and high-concurrency services cannot afford the latency of synchronous blocking system calls. In this lab, students extend Hydra's Linux compatibility subsystem by implementing the \texttt{io\_uring\_setup} and \texttt{io\_uring\_enter} system calls, managing shared submission and completion rings with atomic memory ordering semantics (\texttt{Ordering::Release} and \texttt{Ordering::Acquire}).

\begin{lstlisting}[language=Rust, caption={Submission Queue Processing (\texttt{src/compat/linux/io\_uring\_lab.rs})}]
/// Implements `io_uring_enter` (Syscall #426).
pub fn sys_io_uring_enter(
    ring_fd: i32,
    to_submit: u32,
    min_complete: u32,
    flags: u32,
) -> i64 {
    let task = SCHEDULER.lock().current_task().expect("No task");
    let ring_inst = task.lock().get_io_uring(ring_fd).expect("Bad ring FD");
    let mut ring = ring_inst.lock();

    // 1. Process submitted SQEs
    let submitted = ring.process_sqes();

    // 2. If caller requested blocking wait until min_complete CQEs arrive:
    if min_complete > 0 && ring.cqe_count() < min_complete as usize {
        drop(ring);
        // Suspend current thread on ring completion wait queue
        ring_inst.wait_queue.wait();
    }

    submitted as i64
}
\end{lstlisting}

% ==============================================================================
\section{Lab 8: Mandatory Capability Attenuation and Revocation}
\label{sec:lab8}
% ==============================================================================

\subsection{Architectural Motivation}

In capability-based operating systems, a thread can only perform operations for which it holds an unforgeable capability token. However, a critical design challenge is \textbf{revocation}: if a parent process grants a temporary capability to a child worker, how can the parent revoke that capability later without invalidating the capabilities of unrelated tasks?

In this lab, students implement a \textbf{Capability Derivation Tree (CDT)}. Every time a capability is copied or attenuated (e.g., converting a Read-Write file capability into a Read-Only capability), a node is added to the derivation tree. Revoking a root capability traverses the tree and invalidates all descendant capabilities in $O(N)$ time.

\begin{lstlisting}[language=Rust, caption={Capability Derivation Tree (\texttt{src/capability/cdt.rs})}]
pub struct CapabilityNode {
    pub cap_id: u64,
    pub parent_id: Option<u64>,
    pub rights_mask: u32,
    pub valid: bool,
    pub children: alloc::vec::Vec<u64>,
}

pub struct CapabilityTree {
    nodes: alloc::collections::BTreeMap<u64, CapabilityNode>,
}

impl CapabilityTree {
    /// Recursively revokes a capability and all attenuated children.
    pub fn revoke_cascade(&mut self, root_id: u64) -> usize {
        let mut count = 0;
        let mut stack = alloc::vec![root_id];

        while let Some(current_id) = stack.pop() {
            if let Some(node) = self.nodes.get_mut(&current_id) {
                if node.valid {
                    node.valid = false;
                    count += 1;
                }
                for &child_id in &node.children {
                    stack.push(child_id);
                }
            }
        }

        count
    }
}
\end{lstlisting}

% ==============================================================================
\section{Comprehensive Review and Design Exercises}
\label{sec:all_chapters_exercises}
% ==============================================================================

To conclude our study of modern operating system design and implementation, we present 24 advanced design exercises spanning the entire textbook curriculum:

\begin{enumerate}
    \item \textbf{Microkernel Performance}: Derive the theoretical lower bound on system call latency for an x86\_64 microkernel. What percentage of CPU cycles during an IPC rendezvous are consumed by TLB flushes and cache evictions?
    \item \textbf{Paging Geometry}: Consider an operating system running on a processor with 5-level paging (PML5, 57-bit virtual addresses). Calculate the total virtual address space capacity. How many total page table pages are required to map a single $4\,\text{KiB}$ page located at virtual address \texttt{0x1000\_0000\_0000\_0000}?
    \item \textbf{ZFS Merkle Trees}: In OpenZFS, explain how a silent bit flip occurring in a deep data block at level $L_0$ is detected at the top-level Uberblock during a scrub. Write out the Fletcher-4 recurrence relations.
    \item \textbf{Copy-on-Write Semantics}: When \texttt{fork()} duplicates a process with 1~GiB of allocated virtual memory, how many physical frames are actually allocated by the kernel during the system call? Under what circumstances could an out-of-memory kernel panic occur minutes after a successful fork?
    \item \textbf{I/O Isolation}: Why is the x86\_64 Task State Segment (TSS) I/O Permission Bitmap (IOPB) limited to $64\,\text{KiB}$? How does Hydra ensure that a Ring~3 device driver cannot issue Port I/O instructions to arbitrary legacy hardware?
    \item \textbf{Network Stack Concurrency}: Explain why circular DMA descriptor rings in modern NICs use power-of-two buffer sizes. How does the bitwise mask \texttt{idx \& (N - 1)} avoid division instructions on critical packet reception paths?
    \item \textbf{Wayland Security}: Detail the Confused Deputy problem that arises in X11 when an untrusted application captures keystrokes from a password dialogue box. How does Wayland's private surface protocol eliminate this vulnerability by design?
    \item \textbf{Meltdown and KPTI}: Explain why Kernel Page Table Isolation (KPTI) is necessary to mitigate the Meltdown hardware vulnerability (CVE-2017-5754). What is the exact performance penalty of KPTI on system call intensive workloads?
    \item \textbf{Preemptive Scheduling}: Contrast the behavior of a multi-level feedback queue (MLFQ) scheduler with a proportional-share stride scheduler when handling mixed I/O-bound and CPU-bound workloads.
    \item \textbf{Slab Allocation}: Why does the slab allocator place the \texttt{FreeNode} pointer inside the unused object buffer rather than maintaining an external array of pointers? What cache locality benefit does this provide?
    \item \textbf{Futex Contention}: Why does the Linux kernel provide both \texttt{FUTEX\_WAIT} and \texttt{FUTEX\_WAKE} rather than managing locks purely in user space or purely in supervisor mode?
    \item \textbf{Zero-Copy DMA}: When an Ethernet controller receives a frame via DMA, why must the physical memory buffer be aligned on a 64-byte or 128-byte boundary matching processor cache lines?
    \item \textbf{ZFS Space Maps}: Contrast the Space Map log-structured allocation technique used by OpenZFS with traditional block allocation bit-maps (such as those in ext4).
    \item \textbf{Amdahl's Law in OS Engines}: If 15\% of an operating system kernel's execution time is spent in a serialized global scheduler spinlock, what is the maximum speedup achievable on a 64-core processor?
    \item \textbf{Hardware TLB Invalidation}: When unmapping a virtual memory range spanning 512 continuous pages ($2\,\text{MiB}$), is it more efficient to issue 512 individual \texttt{invlpg} instructions or a single write to \texttt{CR3}? At what threshold does a full TLB reload become preferable?
    \item \textbf{VirtIO Split Rings}: In VirtIO 1.0, why are the Available Ring and Used Ring separated into distinct memory regions? Why does this separation improve performance in multi-vCPU environments?
    \item \textbf{Spectre v2 Mitigations}: How do Retpolines (return trampolines) prevent speculative execution attacks against indirect function pointers in modern C and Rust kernels?
    \item \textbf{System V Calling Convention}: In the System V AMD64 ABI, which register holds the system call number, and which registers hold the first six arguments? Why is \texttt{\%rcx} clobbered by the x86\_64 \texttt{syscall} instruction?
    \item \textbf{Dnode Geometry}: In OpenZFS, how many bonus bytes are available in a 512-byte Dnode for storing inline directory entries (Micro-ZAP) or extended attributes?
    \item \textbf{TCP Sliding Window}: If the TCP receiver window (\texttt{rcv\_wnd}) is 65,535 bytes and the round-trip time (RTT) is $20\,\text{ms}$, calculate the theoretical maximum bandwidth achievable without TCP Window Scaling (RFC 7323).
    \item \textbf{Capability Confinement}: Explain how the confinement problem (Lampson, 1973) is addressed in Hydra's capability model when an untrusted server attempts to leak sensitive data via covert timing channels.
    \item \textbf{Wayland Double-Buffering}: Explain the interaction between \texttt{wl\_surface.attach}, \texttt{wl\_surface.damage}, and \texttt{wl\_surface.commit} during high-frequency window redraws.
    \item \textbf{Synthetic \texttt{/proc} Parsing}: When an application issues \texttt{read()} on \texttt{/proc/self/maps}, why does Hydra synthesize the string on-the-fly rather than reading from a persistent file on disk?
    \item \textbf{Kernel Panic Diagnostics}: In Hydra, when a fatal unhandled exception occurs in Ring~0, what architectural steps does the kernel take to print registers and stack frames to the UEFI GOP linear framebuffer before halting the CPU with \texttt{cli; hlt}?
\end{enumerate}
